% Part: incompleteness
% Chapter: theories-computability
% Section: computably-axiomatizable

\documentclass[../../../include/open-logic-section]{subfiles}

\begin{document}

\olfileid{inc}{tcp}{cax}

\olsection{Teori kang \printtoken{S}{axiomatizable}}

Teori $\Th{T}$ diarani \emph{!!{axiomatizable}} yen nduweni
himpunan aksioma~$A$ kang komputabel. (Nyebut $A$ minangka himpunan
aksioma kanggo $\Th{T}$ tegese $T = \Setabs{!A}{A \Proves !A}$.)
Saben aksiomatisasi bilangan asli kang ``lumrah'' nduweni sipat iki.
Mligi, saben teori kang nduweni himpunan aksioma winates iku
!!{axiomatizable}.

\begin{lem}
Upamana $\Th{T}$ iku !!{axiomatizable}. Mula $\Th{T}$ iku !!{computably
enumerable}.
\end{lem}

\begin{proof}
Upamana $A$ himpunan aksioma kang komputabel kanggo $\Th{T}$. Kanggo
nemtokake apa $!A \in T$, cukup goleka !!a{derivation} saka aksioma
menyang $!A$.

Kanthi tembung liya, $!A$ kalebu ing $\Th{T}$ yen lan mung yen ana
dhaftar winates aksioma $!B_1$, \dots, $!B_k$ ing $A$ lan !!a{derivation}
saka $(!B_1 \land \dots \land !B_k) \lif !A$ ing logika orde siji.
Nanging kita wis ngerti yen saben himpunan kanthi definisi awujud
``ana \dots saengga \dots'' iku !!{c.e.}, angger ``\dots'' kang
kapindho komputabel.
\end{proof}

\end{document}
